\section{The multiplicity boundary floor}
\label{CFStage:BF:BF:section}

This section proves an exact lower bound on the relative control defect of
the theta-generated representatives.  The bound counts every retained
zero with its complete multiplicity.  The source constructions are
A1644, lines 2576--2718, and A1771, lines 4477--4835, in the authoritative
\texttt{audit\_segment\_U0041\_U0054.md}.  The rank-two identity itself is
A1771, lines 4688--4725.  The invariant trace and the invariant-subspace
argument below are new deductions from those source constructions.

All matrix coordinates on the packet algebra are the original polynomial
coordinates.  The generator is \(D=-x\partial_x\), the arithmetic
multiplier is \(g=2\xi\), and the Hilbert-space measure is \(dx\) on the
positive half-line.  The local extension unit, the unscaled orthogonal
vectors, their norm squares, and all nilpotent jets are retained.

\subsection{The original source, packet algebra, and arithmetic section}
\label{CFStage:BF:BF:source}

Write
\[
\xi(s)=\frac12 s(s-1)\pi^{-s/2}\Gamma(s/2)\zeta(s),
\qquad g(s)=2\xi(s),
\]
with the entire continuations understood at the apparent singularities.
Let
\begin{equation}
\begin{split}
V={}&\left\{\phi\in\mathcal S(\mathbb R):
  \phi\text{ is even},\ \phi(0)=0,\
  \int_{\mathbb R}\phi(x)\,dx=0\right\},\\
\mathcal B={}&\left\{F\in C^\infty(\mathbb R_{>0}):
  \sup_{x>0}x^b|D^jF(x)|<\infty\
  \text{for all }b\in\mathbb Z,\ j\ge0\right\},\\
D={}&-x\partial_x,\qquad
(\Theta\phi)(x)=\sum_{a\ne0}\phi(ax)
               =2\sum_{a\ge1}\phi(ax).
\end{split}
\tag{BF.01}\label{CFStage:BF:BF:eq:spaces}
\end{equation}
Here \(\mathcal S(\mathbb R)\) is the Schwartz space, and the summation
indices \(a\) in \(\Theta\) are integers.  The symbol \(\mathcal B\)
denotes exactly the strong-Schwartz target specified in A1694,
lines 3049--3073.  Put
\[
Q=\mathcal B/\Theta V,\qquad q:\mathcal B\longrightarrow Q,
\qquad
\mathcal H_+=L^2(\mathbb R_{>0},dx).
\]
The scalar product convention is
\[
\langle F,H\rangle_+
 =\int_0^\infty\overline{F(x)}H(x)\,dx,
\]
antilinear in the first argument.  Adjoint signs on maps into
\(\mathcal H_+\) use this scalar product and the standard scalar product
in the stated polynomial coordinates on their finite-dimensional
domains.

Retain the actual source function and its derivative sequence:
\begin{equation}
\begin{split}
\phi_*(x)&=(4\pi^2x^4-6\pi x^2)e^{-\pi x^2},\\
f_0&=\Theta\phi_*,
\qquad \phi_j=D^j\phi_*,
\qquad f_j=\Theta\phi_j=D^jf_0\quad(j\ge0),\\
(\mathcal MF)(s)&=\int_0^\infty F(x)x^s\,\frac{dx}{x},
\qquad \mathcal Mf_0=g.
\end{split}
\tag{BF.02}\label{CFStage:BF:BF:eq:theta-source}
\end{equation}
The source identity \(\mathcal M(D^jf_0)(s)=s^jg(s)\), with these exact
constants, is A1644, lines 2566--2574, and A1771, lines 4479--4529.
The commutation \(D\Theta\phi=\Theta D\phi\) uses the same generator on
the source and target.  Differentiation of the defining series is
justified on each compact subinterval of \((0,\infty)\) by the
Schwartz estimates.

Let \(Z\) be a nonempty finite set of actual zeros of \(g\).  For
\(\rho\in Z\), let \(m_\rho\ge1\) be its complete order of vanishing.
Define
\begin{equation}
h_Z(t)=\prod_{\rho\in Z}(t-\rho)^{m_\rho},\qquad
d=\deg h_Z=\sum_{\rho\in Z}m_\rho,\qquad
E_Z=\mathbb C[t]/(h_Z),\qquad
A=M_t:E_Z\longrightarrow E_Z.
\tag{BF.03}\label{CFStage:BF:BF:eq:packet}
\end{equation}
The ordered coordinates on \(E_Z\) are the classes of
\(1,t,\ldots,t^{d-1}\).  Multiplication by \(t\) is not replaced by a
shifted generator.  The empty packet is treated separately at the end.

For a function holomorphic near every member of \(Z\), \(j_{h_Z}\)
means its complete collection of jets of orders \(0,\ldots,m_\rho-1\),
identified with \(E_Z\) by the Chinese remainder isomorphism.
In particular, the finite Mellin-jet map is
\[
J_Z:\mathcal B\longrightarrow E_Z,\qquad
J_ZF=j_{h_Z}(\mathcal MF).
\]
In the local coordinate \(z=s-\rho\), the \(\rho\)-component is
\[
\sum_{j=0}^{m_\rho-1}
\frac{(\mathcal MF)^{(j)}(\rho)}{j!}\,z^j
       \pmod{z^{m_\rho}}.
\]
The map retains the factorials and the full jet length.

The actual extension unit is
\begin{equation}
\varepsilon_Z
 =\left(j_{h_Z}(g/h_Z)\right)^{-1}
 =j_{h_Z}(h_Z/g)\in E_Z^\times.
\tag{BF.04}\label{CFStage:BF:BF:eq:extension-unit}
\end{equation}
To see why it is a unit, write
\[
g(\rho+z)=z^{m_\rho}w_\rho(z),\qquad w_\rho(0)\ne0.
\]
The polynomial \(h_Z(\rho+z)\) has precisely the same zero order.
Thus \(g/h_Z\) and \(h_Z/g\) are mutually inverse holomorphic units
near every selected zero.  Equation \eqref{CFStage:BF:BF:eq:extension-unit}
keeps their complete jets.  This finite-algebra unit is distinct from
the supported-zero scalar \(e\) of the original split-support
construction.

For \(u\in E_Z\), let \(R_Z(u)\) be the unique polynomial of degree
less than \(d\) representing \(\varepsilon_Zu\), and define the
extension row
\begin{equation}
\ell_Z:E_Z\longrightarrow\mathbb C,\qquad
\ell_Z(u)=[t^{d-1}]R_Z(u).
\tag{BF.05}\label{CFStage:BF:BF:eq:extension-row}
\end{equation}
The established arithmetic section in A1644, lines 2594--2718, and
A1771, lines 4499--4510, is the map
\[
s_Z:E_Z\longrightarrow\mathcal B
\]
with the identities
\begin{equation}
J_Zs_Z=I_{E_Z},\qquad J_Z\Theta=0,\qquad
qs_Z=\sigma_Z:E_Z\hookrightarrow Q,\qquad
Ds_Z-s_ZA=f_0\ell_Z.
\tag{BF.06}\label{CFStage:BF:BF:eq:section}
\end{equation}
These are the inherited source identities used below, with the actual
source section and its analytic construction.  For a one-point packet
at a simple zero, their coefficient is exactly
\[
\ell_Z(1)=\frac{1}{g'(\rho)}
         =\frac{1}{2\xi'(\rho)}.
\]
No replacement of this coefficient by a unit-free scalar is made.

\subsection{Theta relation spaces and the exact rank-two boundary}
\label{CFStage:BF:BF:rank-two}

Define
\begin{equation}
V_n=\operatorname{span}_{\mathbb C}\{\phi_0,\ldots,\phi_n\},
\qquad
L_n=\Theta V_n
   =\operatorname{span}_{\mathbb C}\{f_0,\ldots,f_n\},
\qquad n\ge0.
\tag{BF.07}\label{CFStage:BF:BF:eq:relation-spaces}
\end{equation}
The vectors \(f_j\) are linearly independent.  Indeed, a finite
dependence has Mellin transform \(g(s)P(s)=0\), where
\(P(s)=\sum_jc_js^j\).  Since \(g\) is not identically zero, it is
nonzero on a nonempty open set.  Therefore \(P\) vanishes on that
open set and is the zero polynomial.  Every coefficient \(c_j\)
is zero.

Let \(F_n:\mathbb C^{n+1}\to\mathcal H_+\) have columns
\(f_0,\ldots,f_n\), and set
\begin{equation}
H_n=F_n^*F_n,\qquad
P_n=F_nH_n^{-1}F_n^*,\qquad
R_n=(I-P_n)s_Z,\qquad
G_n=R_n^*R_n.
\tag{BF.08}\label{CFStage:BF:BF:eq:representatives}
\end{equation}
Independence gives \(H_n\succ0\).  Direct multiplication gives
\(P_n^2=P_n=P_n^*\), and its range is \(L_n\); it is the actual
orthogonal projection onto the specified finite theta relation space.
The expression
\[
b_n(u)=\sum_{j=0}^n
\bigl(H_n^{-1}F_n^*s_Zu\bigr)_j\phi_j
\]
defines a map \(b_n:E_Z\to V_n\) satisfying \(P_ns_Z=\Theta b_n\).
Consequently, by \eqref{CFStage:BF:BF:eq:section},
\begin{equation}
qR_n=\sigma_Z,\qquad J_ZR_n=I_{E_Z},\qquad
G_n=s_Z^*s_Z-s_Z^*F_nH_n^{-1}F_n^*s_Z\succ0.
\tag{BF.09}\label{CFStage:BF:BF:eq:exact-class}
\end{equation}
For the strict inequality, \(v^*G_nv=\|R_nv\|_+^2=0\) implies
\(R_nv=0\) and hence \(v=J_ZR_nv=0\).  Thus all inverses of \(G_n\)
used below are inverses of strictly positive matrices at the stated
finite level.

The original theta boundary and the associated control matrix are
\begin{equation}
\begin{split}
B_n&=DR_n-R_nA=\Theta K_n,\\
K_n&=\phi_*\ell_Z-D_Vb_n+b_nA,\\
W_n&=A^*G_n+G_nA-G_n.
\end{split}
\tag{BF.10}\label{CFStage:BF:BF:eq:defect}
\end{equation}
Here \(D_V\) is \(D\) on \(V\).  The boundary formula follows by
substituting \(R_n=s_Z-\Theta b_n\) into
\eqref{CFStage:BF:BF:eq:section}.  It uses only original theta boundaries.

For \(F,H\in\mathcal B\), integration by parts gives
\begin{equation}
\langle DF,H\rangle_++\langle F,DH\rangle_+
=-\left[x\overline{F(x)}H(x)\right]_{0}^{\infty}
 +\langle F,H\rangle_+
=\langle F,H\rangle_+.
\tag{BF.11}\label{CFStage:BF:BF:eq:parts}
\end{equation}
The defining estimates of \(\mathcal B\) make the boundary expression
zero at both endpoints.  Applying \eqref{CFStage:BF:BF:eq:parts} to columns of
\(R_n\), and substituting \(DR_n=R_nA+B_n\), proves
\begin{equation}
W_n=-(R_n^*B_n+B_n^*R_n).
\tag{BF.12}\label{CFStage:BF:BF:eq:cross-defect}
\end{equation}
The coefficient of \(G_n\) in \eqref{CFStage:BF:BF:eq:defect} is exactly one
because the measure in \eqref{CFStage:BF:BF:eq:parts} is \(dx\).
Even extension to the full line has squared norm twice this norm;
it is not silently substituted here.

Construct the unscaled orthogonal vectors
\begin{equation}
u_0=f_0,\qquad
u_n=(I-P_{n-1})f_n\quad(n\ge1),\qquad
\mathfrak h_n=\langle u_n,u_n\rangle_+>0,\qquad
r_n=\mathfrak h_n^{-1}u_n^*s_Z:E_Z\to\mathbb C.
\tag{BF.13}\label{CFStage:BF:BF:eq:unscaled-rows}
\end{equation}
The operator \(u_n^*:\mathcal H_+\to\mathbb C\) sends \(F\) to
\(\langle u_n,F\rangle_+\).  Each \(u_n\) has leading coefficient
one on \(f_n\); none is divided by its norm.  Since the \(f_j\) are
real functions, their Gram matrices are real, and each \(u_n\) is
real.

Degree and orthogonality determine the recurrence
\begin{equation}
\begin{split}
Du_0&=u_1+\frac12u_0,\\
Du_n&=u_{n+1}+\frac12u_n
          -\frac{\mathfrak h_n}{\mathfrak h_{n-1}}u_{n-1}
          \qquad(n\ge1).
\end{split}
\tag{BF.14}\label{CFStage:BF:BF:eq:unscaled-recurrence}
\end{equation}
Here is the complete coefficient calculation.  As
\(u_n=f_n+\sum_{j<n}c_jf_j\), \(Du_n\) belongs to \(L_{n+1}\)
and has coefficient one on \(u_{n+1}\).
For \(j\le n-2\), \eqref{CFStage:BF:BF:eq:parts} gives
\[
\langle u_j,Du_n\rangle_+
=\langle u_j,u_n\rangle_+-\langle Du_j,u_n\rangle_+=0,
\]
because \(Du_j\in L_{j+1}\subseteq L_{n-1}\).
Reality and \eqref{CFStage:BF:BF:eq:parts} give
\[
2\langle u_n,Du_n\rangle_+=\mathfrak h_n,
\]
so the full diagonal coefficient is \(1/2\), not merely its real
part.  For \(n\ge1\),
\[
\langle u_{n-1},Du_n\rangle_+
=-\langle Du_{n-1},u_n\rangle_+=-\mathfrak h_n.
\]
Dividing only this coefficient by the corresponding original norm
square \(\mathfrak h_{n-1}\) yields the last term of
\eqref{CFStage:BF:BF:eq:unscaled-recurrence}.

The section coefficients satisfy
\begin{equation}
P_ns_Z=\sum_{j=0}^nu_jr_j,\qquad
\frac{\mathfrak h_{n+1}}{\mathfrak h_n}r_{n+1}
=r_n\left(\frac12I-A\right)+r_{n-1}
   -\mathbf 1_{\{n=0\}}\ell_Z,\qquad r_{-1}=0.
\tag{BF.15}\label{CFStage:BF:BF:eq:row-recurrence}
\end{equation}
The indicator \(\mathbf 1_{\{n=0\}}\) equals one for \(n=0\)
and zero otherwise.  To check the row equation, pair
\(Ds_Z=s_ZA+u_0\ell_Z\) against \(u_n\).
The right side is
\[
\mathfrak h_nr_nA+\mathbf 1_{\{n=0\}}\mathfrak h_0\ell_Z.
\]
By \eqref{CFStage:BF:BF:eq:parts} and \eqref{CFStage:BF:BF:eq:unscaled-recurrence}, the
left side is
\[
\frac12\mathfrak h_nr_n-\mathfrak h_{n+1}r_{n+1}
+\mathfrak h_nr_{n-1},
\]
where the last term is zero for \(n=0\).  Rearrangement gives
\eqref{CFStage:BF:BF:eq:row-recurrence}.  The initial extension functional is
therefore fixed by the original arithmetic section.

We now obtain the source's rank-two formula with its exact sign.
Since
\[
B_n=f_0\ell_Z-DP_ns_Z+P_ns_ZA,
\]
the first and third terms have image in \(L_n\), orthogonal to
the image of \(R_n\).  Equation \eqref{CFStage:BF:BF:eq:cross-defect} becomes
\[
W_n=R_n^*DP_ns_Z+(DP_ns_Z)^*R_n.
\]
Equations \eqref{CFStage:BF:BF:eq:unscaled-recurrence} and
\eqref{CFStage:BF:BF:eq:row-recurrence} give
\[
(I-P_n)DP_ns_Z=u_{n+1}r_n,\qquad
u_{n+1}^*R_n=u_{n+1}^*s_Z
             =\mathfrak h_{n+1}r_{n+1}.
\]
It follows that
\begin{equation}
\boxed{\;
W_n=\mathfrak h_{n+1}
       \left(r_{n+1}^*r_n+r_n^*r_{n+1}\right).
\;}
\tag{BF.16}\label{CFStage:BF:BF:eq:rank-two}
\end{equation}
Also \(R_{n+1}=R_n-u_{n+1}r_{n+1}\), with \(R_{n+1}\)
orthogonal to \(u_{n+1}\).  Taking Gram matrices proves the exact
decrement
\begin{equation}
G_n-G_{n+1}
 =\mathfrak h_{n+1}r_{n+1}^*r_{n+1}\succeq0.
\tag{BF.17}\label{CFStage:BF:BF:eq:gram-decrement}
\end{equation}
These formulas are A1771, lines 4688--4734.  The rank bound in
\eqref{CFStage:BF:BF:eq:rank-two} concerns the boundary control form on \(E_Z\);
the dimension \(d\), and the full operator \(A\), have not been
reduced to two dimensions.

For clarity, the retained source-support transition is also exact.
Write \(Q_{V_n}=\mathcal B/\Theta V_n=\mathcal B/L_n\).
Since \(D\Theta V_n\subseteq\Theta V_{n+1}\), differentiation gives
a well-defined map
\[
\overline D_n:Q_{V_n}\longrightarrow Q_{V_{n+1}},
\qquad [F]\longmapsto[DF].
\]
Before that transition, the boundary class is
\[
[B_nv]_{Q_{V_n}}=-[u_{n+1}]_{Q_{V_n}}r_n(v).
\]
The layer map
\[
\Theta V_{n+1}/\Theta V_n
\longrightarrow L_{n+1}\cap L_n^\perp,\qquad
[z]\longmapsto(I-P_n)z
\]
is well-defined and injective because its kernel in \(L_{n+1}\)
is \(L_n\).  It is surjective because every vector in the displayed
orthogonal complement is fixed by \(I-P_n\).
The subsequent support transport makes this particular class the
supported zero in \(Q_{V_{n+1}}\).  Its earlier pairing in
\eqref{CFStage:BF:BF:eq:rank-two} is retained.  These are the actual maps of
A1771, lines 4736--4773.

\subsection{Relative control spectrum and its invariant trace}
\label{CFStage:BF:BF:trace}

For \(G_n\succ0\), define its least two-sided control excess by
\begin{equation}
\epsilon_n=\inf\{\epsilon\ge0:
               -\epsilon G_n\preceq W_n\preceq\epsilon G_n\}.
\tag{BF.18}\label{CFStage:BF:BF:eq:excess-definition}
\end{equation}
Here \(X\preceq Y\) means that \(Y-X\) is a positive semidefinite
Hermitian matrix.  The operator \(G_n^{-1}W_n\) is self-adjoint in
the positive inner product \(\langle u,v\rangle_{G_n}=u^*G_nv\):
\[
\langle u,G_n^{-1}W_nv\rangle_{G_n}=u^*W_nv
=\langle G_n^{-1}W_nu,v\rangle_{G_n}.
\]
Thus \(\epsilon_n\) is the largest absolute value of its eigenvalues.
This uses the actual metric and operator and makes no assumption that
\(A\) is normal for this metric.

In the unchanged coordinates, set
\begin{equation}
\mathsf a_n=r_nG_n^{-1}r_n^*,\qquad
\mathsf b_n=r_{n+1}G_n^{-1}r_{n+1}^*,\qquad
\mathsf c_n=r_nG_n^{-1}r_{n+1}^*.
\tag{BF.19}\label{CFStage:BF:BF:eq:abc}
\end{equation}
The exact maps
\begin{equation}
\begin{split}
\mathcal T_n&:E_Z\longrightarrow\mathbb C^2,
\qquad\mathcal T_nv=(r_nv,r_{n+1}v),\\
\mathcal U_n&:\mathbb C^2\longrightarrow E_Z,\\
\mathcal U_n(a,b)
 &=\mathfrak h_{n+1}
    \left(G_n^{-1}r_{n+1}^*a+G_n^{-1}r_n^*b\right)
\end{split}
\tag{BF.20}\label{CFStage:BF:BF:eq:factor-maps}
\end{equation}
satisfy
\[
\mathcal U_n\mathcal T_n=G_n^{-1}W_n,\qquad
\mathcal T_n\mathcal U_n
=\mathfrak h_{n+1}
  \begin{pmatrix}
   \mathsf c_n&\mathsf a_n\\
   \mathsf b_n&\overline{\mathsf c_n}
  \end{pmatrix}.
\]
For every nonzero eigenvalue, applying \(\mathcal T_n\) and
\(\mathcal U_n\) gives mutually inverse eigenspace maps up to that
nonzero scalar.  Therefore the nonzero eigenvalues of the two
composites agree.  Cauchy--Schwarz for the positive form on row
covectors induced by \(G_n^{-1}\) gives
\(\mathsf a_n\mathsf b_n\ge|\mathsf c_n|^2\).
The two roots of the displayed \(2\)-by-\(2\) characteristic
polynomial are consequently
\begin{equation}
\lambda_{n,\pm}
=\mathfrak h_{n+1}
 \left(\operatorname{Re}\mathsf c_n
 \ \pm\
 \sqrt{\mathsf a_n\mathsf b_n-
               (\operatorname{Im}\mathsf c_n)^2}\right),
\qquad \lambda_{n,+}\ge0\ge\lambda_{n,-}.
\tag{BF.21}\label{CFStage:BF:BF:eq:two-roots}
\end{equation}
Zero roots are included in this notation even when the corresponding
eigenspace does not occur in \(E_Z\).
There are no other nonzero eigenvalues by the factorization
\eqref{CFStage:BF:BF:eq:factor-maps}.  Thus
\begin{equation}
\epsilon_n=\mathfrak h_{n+1}
\left(
 |\operatorname{Re}\mathsf c_n|
 +\sqrt{\mathsf a_n\mathsf b_n-
                  (\operatorname{Im}\mathsf c_n)^2}
\right).
\tag{BF.22}\label{CFStage:BF:BF:eq:explicit-excess}
\end{equation}
This is the exact expression in A1771, lines 4775--4821.
In particular, \(G_n^{-1}\) remains present in every relative
coefficient.

The trace of this control operator is independent of the chosen
positive representative metric.  Cyclicity of finite-dimensional
trace and \eqref{CFStage:BF:BF:eq:defect} give
\begin{equation}
\begin{split}
\operatorname{tr}(G_n^{-1}W_n)
&=\operatorname{tr}(G_n^{-1}A^*G_n)
  +\operatorname{tr}A-d\\
&=2\operatorname{Re}\operatorname{tr}A-d\\
&=\sum_{\rho\in Z}m_\rho(2\operatorname{Re}\rho-1).
\end{split}
\tag{BF.23}\label{CFStage:BF:BF:eq:trace-invariant}
\end{equation}
For the final equality, the Chinese remainder decomposition preserves
the local coordinates \(z=t-\rho\).  On the length-\(m_\rho\)
summand, multiplication by \(t\) is \(\rho I+M_z\).  In the ordered
basis \(1,z,\ldots,z^{m_\rho-1}\), \(M_z\) has zero diagonal and
trace zero.  Thus its full trace contribution is \(m_\rho\rho\).
No generalized eigenvector has been deleted.

Taking the trace directly in \eqref{CFStage:BF:BF:eq:rank-two}, and using
\eqref{CFStage:BF:BF:eq:abc}, yields the exact row constraint
\begin{equation}
\boxed{\;
2\mathfrak h_{n+1}\operatorname{Re}\mathsf c_n
 =\sum_{\rho\in Z}m_\rho(2\operatorname{Re}\rho-1).
\;}
\tag{BF.24}\label{CFStage:BF:BF:eq:row-trace}
\end{equation}

Define the reflection \(\iota\rho=1-\bar\rho\).
A reflection-stable packet means that \(\iota Z=Z\) and that
\(m_{\iota\rho}=m_\rho\) at every member.
For such a packet the contributions of every two-point orbit in
\eqref{CFStage:BF:BF:eq:trace-invariant} cancel.  At a fixed point of \(\iota\),
\(\operatorname{Re}\rho=1/2\), so its contribution is zero.
Since \(\mathfrak h_{n+1}>0\),
\begin{equation}
\begin{split}
\operatorname{tr}(G_n^{-1}W_n)&=0,
\qquad\operatorname{Re}\mathsf c_n=0,\\
\lambda_{n,+}&=\epsilon_n=-\lambda_{n,-},\\
\epsilon_n
 &=\mathfrak h_{n+1}
   \sqrt{\mathsf a_n\mathsf b_n-
                 (\operatorname{Im}\mathsf c_n)^2}.
\end{split}
\tag{BF.25}\label{CFStage:BF:BF:eq:reflection-trace}
\end{equation}
If \(W_n\ne0\), the two nonzero eigenvalues are
\(\epsilon_n\) and \(-\epsilon_n\), each once.
If \(\epsilon_n=0\), the self-adjoint operator \(G_n^{-1}W_n\)
has only the zero eigenvalue and hence is zero.  Trace cancellation
alone does not set the matrix \(W_n\) to zero.

\subsection{The invariant-subspace projection argument}
\label{CFStage:BF:BF:projection-proof}

The next estimates hold for an arbitrary actual packet \(Z\);
reflection stability will be used only for their final combination.
Define
\[
Z_+=\{\rho\in Z:\operatorname{Re}\rho>1/2\},\qquad
Z_-=\{\rho\in Z:\operatorname{Re}\rho<1/2\}.
\]
Let
\begin{equation}
E_+=\bigoplus_{\rho\in Z_+}\mathbb C[z]/(z^{m_\rho}),
\qquad
E_-=\bigoplus_{\rho\in Z_-}\mathbb C[z]/(z^{m_\rho}).
\tag{BF.26}\label{CFStage:BF:BF:eq:spectral-subspaces}
\end{equation}
The inclusions \(I_\pm:E_\pm\to E_Z\) are the actual Chinese
remainder inclusions, not coordinate projections chosen to make
\(A\) normal.  Explicitly, put \(h_\rho=(t-\rho)^{m_\rho}\) and
\(H_\rho=h_Z/h_\rho\).  Since \(H_\rho\) and \(h_\rho\) are coprime,
choose \(b_\rho(t)\) with \(H_\rho b_\rho\equiv1\pmod{h_\rho}\),
and let \(e_\rho\) be the class of \(H_\rho b_\rho\) in \(E_Z\).
It has residue one at the \(\rho\)-summand and zero at all other
summands.  The maps in \eqref{CFStage:BF:BF:eq:spectral-subspaces} are
\[
I_\pm((u_\rho)_\rho)
 =\sum_{\rho\in Z_\pm}e_\rho\,u_\rho(t-\rho).
\]
Changing the representative \(b_\rho\) changes nothing in \(E_Z\).
All original polynomial coordinates and every full local length
remain in these maps.

Let \(A_\pm\) be multiplication by \(\rho+z\) on each local
summand of \(E_\pm\).  The intertwining identities are
\begin{equation}
AI_\pm=I_\pm A_\pm.
\tag{BF.27}\label{CFStage:BF:BF:eq:invariant-inclusions}
\end{equation}
In particular the images are invariant generalized eigenspaces, with
their nilpotent parts unchanged.

Fix \(n\ge0\).  For a nonzero \(E_\pm\), define
\begin{equation}
G_\pm=I_\pm^*G_nI_\pm,\qquad
P_\pm=I_\pm G_\pm^{-1}I_\pm^*G_n.
\tag{BF.28}\label{CFStage:BF:BF:eq:orthogonal-projections}
\end{equation}
The notation \(G_\pm,P_\pm\) in this subsection refers to this
fixed value of \(n\).
Since \(I_\pm\) is injective and \(G_n\succ0\), \(G_\pm\succ0\).
Multiplication gives \(P_\pm^2=P_\pm\),
\(P_\pm I_\pm=I_\pm\), and
\[
P_\pm^*G_n=G_nP_\pm.
\]
Thus \(P_\pm\) is the \(G_n\)-orthogonal projection onto
\(\operatorname{im}I_\pm\).  When \(E_\pm=\{0\}\), define
\(P_\pm=0\) directly; no inverse of a nonpositive-dimensional
Gram matrix is needed.

For nonzero \(E_\pm\), the restricted defect is exact.
By \eqref{CFStage:BF:BF:eq:invariant-inclusions},
\begin{equation}
I_\pm^*W_nI_\pm
 =A_\pm^*G_\pm+G_\pm A_\pm-G_\pm.
\tag{BF.29}\label{CFStage:BF:BF:eq:restricted-defect}
\end{equation}
For nonzero \(E_\pm\), cyclicity of the trace of the rectangular
composites, followed by \eqref{CFStage:BF:BF:eq:restricted-defect}, gives
\begin{equation}
\begin{split}
\operatorname{tr}(P_\pm G_n^{-1}W_n)
&=\operatorname{tr}\!\left(
  G_\pm^{-1}I_\pm^*W_nI_\pm\right)\\
&=2\operatorname{Re}\operatorname{tr}A_\pm-\dim E_\pm\\
&=2\sum_{\rho\in Z_\pm}
        m_\rho(\operatorname{Re}\rho-1/2).
\end{split}
\tag{BF.30}\label{CFStage:BF:BF:eq:compressed-trace}
\end{equation}
For \(E_\pm=\{0\}\), the first and last expressions in
\eqref{CFStage:BF:BF:eq:compressed-trace} are both zero and remain equal;
the inverse-containing expression is not invoked.
The last trace evaluation again uses
\(\rho I+M_z\) on each full local summand.

The rank-two factorization gives a stronger estimate for these
compressed traces than a generic rank bound would give.
Indeed, the Hermitian form \(v^*W_nv\) is the pullback under
\(v\mapsto(r_nv,r_{n+1}v)\) of
\begin{equation}
\mathfrak h_{n+1}
\begin{pmatrix}0&1\\1&0\end{pmatrix}.
\tag{BF.31}\label{CFStage:BF:BF:eq:exchange-form}
\end{equation}
This form has one positive and one negative direction.
To verify the inherited inertia bound directly, a positive definite
subspace for \(W_n\) must map injectively under
\(v\mapsto(r_nv,r_{n+1}v)\); otherwise its kernel would contain
a nonzero vector on which the form is zero.
Its image is a positive definite subspace of
\eqref{CFStage:BF:BF:eq:exchange-form}, so has dimension at most one.
The same reasoning applies to a negative definite subspace.
Therefore \(G_n^{-1}W_n\), self-adjoint in the \(G_n\) metric,
has at most one strictly positive and at most one strictly negative
eigenvalue.

Use \(\lambda_{n,+}\) and \(\lambda_{n,-}\) from
\eqref{CFStage:BF:BF:eq:two-roots}.  Let \(\Pi_+\) and \(\Pi_-\)
be their \(G_n\)-orthogonal spectral projections when the
corresponding eigenvalue is strictly positive or strictly negative.
When a sign is absent, set both its eigenvalue and its projection
equal to zero.  The finite-dimensional spectral theorem gives
\begin{equation}
G_n^{-1}W_n
 =\lambda_{n,+}\Pi_++\lambda_{n,-}\Pi_-,
\qquad \operatorname{rank}\Pi_\pm\le1.
\tag{BF.32}\label{CFStage:BF:BF:eq:spectral-projections}
\end{equation}
All remaining spectral directions have eigenvalue zero.

Let \(P\) be any \(G_n\)-orthogonal projection.
For either of the nonzero projections \(\Pi_\pm\),
\(\Pi_\pm P\Pi_\pm\) acts on a one-dimensional subspace as a
positive contraction.  Positivity follows because
\[
\langle x,\Pi_\pm P\Pi_\pm x\rangle_{G_n}
=\langle P\Pi_\pm x,P\Pi_\pm x\rangle_{G_n}\ge0,
\]
and the upper bound follows from
\(\|P\Pi_\pm x\|_{G_n}\le\|\Pi_\pm x\|_{G_n}\).
Its eigenvalue belongs to \([0,1]\).
Cyclicity and \(\Pi_\pm^2=\Pi_\pm\) then prove
\begin{equation}
0\le\operatorname{tr}(P\Pi_\pm)
 =\operatorname{tr}(\Pi_\pm P\Pi_\pm)\le1.
\tag{BF.33}\label{CFStage:BF:BF:eq:projection-overlap}
\end{equation}
The equality and both inequalities also hold when \(\Pi_\pm=0\).
Combining \eqref{CFStage:BF:BF:eq:spectral-projections} and
\eqref{CFStage:BF:BF:eq:projection-overlap} gives
\begin{equation}
\lambda_{n,-}
\le\operatorname{tr}(P G_n^{-1}W_n)
\le\lambda_{n,+}.
\tag{BF.34}\label{CFStage:BF:BF:eq:trace-sandwich}
\end{equation}
For example the upper bound is
\(\lambda_{n,+}\operatorname{tr}(P\Pi_+)
+\lambda_{n,-}\operatorname{tr}(P\Pi_-)\le\lambda_{n,+}\);
the lower bound follows with the two signs interchanged.

Apply \eqref{CFStage:BF:BF:eq:trace-sandwich} to the actual projections
\(P_+\) and \(P_-\) in \eqref{CFStage:BF:BF:eq:compressed-trace}.
For every packet and every \(n\ge0\), this proves
\begin{equation}
\boxed{\begin{aligned}
\lambda_{n,+}
 &\ge2\sum_{\rho\in Z_+}
          m_\rho(\operatorname{Re}\rho-1/2),\\
-\lambda_{n,-}
 &\ge2\sum_{\rho\in Z_-}
          m_\rho(1/2-\operatorname{Re}\rho).
\end{aligned}}
\tag{BF.35}\label{CFStage:BF:BF:eq:one-sided-mass}
\end{equation}
These estimates do not require reflection stability or
semisimplicity.

For a reflection-stable packet, reflection bijects \(Z_+\) and
\(Z_-\), preserves multiplicity, and negates
\(\operatorname{Re}\rho-1/2\).  Thus the two positive distance
sums in \eqref{CFStage:BF:BF:eq:one-sided-mass} are equal.
Equation \eqref{CFStage:BF:BF:eq:reflection-trace} identifies the two control
eigenvalues with \(\epsilon_n\) and \(-\epsilon_n\).
Their combination is the claimed multiplicity boundary floor:
\begin{equation}
\boxed{\begin{aligned}
\epsilon_n
 &=\mathfrak h_{n+1}
   \sqrt{\mathsf a_n\mathsf b_n-
                 (\operatorname{Im}\mathsf c_n)^2}\\
 &\ge\sum_{\rho\in Z}m_\rho
               \left|\operatorname{Re}\rho-\frac12\right|
 \qquad(n\ge0).
\end{aligned}}
\tag{BF.36}\label{CFStage:BF:BF:eq:boundary-floor}
\end{equation}
The restriction traces in \eqref{CFStage:BF:BF:eq:compressed-trace} show
explicitly why complete multiplicities occur.  The proof includes
every generalized eigenspace and uses its induced metric.  No
normality of \(A\), deletion of nilpotent jets, or positivity of
the arithmetic Weil form has entered.

\subsection{Exact consequences for a putative off-line zero}
\label{CFStage:BF:BF:counterfactual}

An individual eigenvector gives a useful comparison.
If \(0\ne v\in E_Z\) satisfies \(Av=\rho v\), then directly from
\eqref{CFStage:BF:BF:eq:defect},
\begin{equation}
v^*W_nv=(2\operatorname{Re}\rho-1)v^*G_nv.
\tag{BF.37}\label{CFStage:BF:BF:eq:eigenline}
\end{equation}
Such a vector exists in every local summand: the class
\(z^{m_\rho-1}\), embedded by \(e_\rho\), is nonzero and is
annihilated by \(A-\rho I\).  By
\eqref{CFStage:BF:BF:eq:excess-definition}, \eqref{CFStage:BF:BF:eq:eigenline} gives
\(\epsilon_n\ge|2\operatorname{Re}\rho-1|\).
Equation \eqref{CFStage:BF:BF:eq:boundary-floor} includes the additional
constraints from all other roots and their multiplicities.

For a reflected off-line pair
\(\{\rho,1-\bar\rho\}\), with complete multiplicity \(m\) at both
points, put \(\delta=\operatorname{Re}\rho-1/2\ne0\).
The two points are distinct, and their contributions to the sum
in \eqref{CFStage:BF:BF:eq:boundary-floor} give
\begin{equation}
\epsilon_n\ge2m|\delta|
 =m|2\operatorname{Re}\rho-1|
 \qquad(n\ge0).
\tag{BF.38}\label{CFStage:BF:BF:eq:pair}
\end{equation}
The real structure and functional equation of the retained
xi-function are \(g(\bar s)=\overline{g(s)}\) and \(g(1-s)=g(s)\).
Composition of a local Taylor expansion with conjugation or the
affine map \(s\mapsto1-s\) preserves its first nonzero degree, so
these symmetries preserve the complete vanishing multiplicities.
Consequently, when the four points
\(\rho,\bar\rho,1-\rho,1-\bar\rho\) are distinct, they form a
reflection-stable packet with the same complete multiplicity \(m\).
Their exact bound is
\begin{equation}
\epsilon_n\ge4m|\delta|
 =2m|2\operatorname{Re}\rho-1|
 \qquad(n\ge0).
\tag{BF.39}\label{CFStage:BF:BF:eq:quartet}
\end{equation}
For a collapsed symmetry orbit, only the actual distinct points
with their actual multiplicities are counted.  Formula
\eqref{CFStage:BF:BF:eq:pair} already treats the two-point reflection orbit.

In particular, a reflection-stable actual packet containing an
off-line root forces a strictly positive right side of
\eqref{CFStage:BF:BF:eq:boundary-floor} at every finite level.
Then \(W_n\) is nonzero, indefinite, and of rank exactly two,
with generalized eigenvalues \(\epsilon_n,-\epsilon_n\).
The strict positivity follows from the packet sum; both signs
and exact rank follow from \eqref{CFStage:BF:BF:eq:reflection-trace}
and \eqref{CFStage:BF:BF:eq:rank-two}.  For an arbitrary packet without
reflection stability, the valid statements are instead the
separate inequalities \eqref{CFStage:BF:BF:eq:one-sided-mass}; its trace
need not vanish.

This lower bound is on the relative expression containing
\(G_n^{-1}\).  The absolute monotonicity
\eqref{CFStage:BF:BF:eq:gram-decrement} does not supply an upper bound on
that inverse-weighted expression.
The homotopy freedom in A1771, lines 4432--4475, does not alter
this conclusion.  Retain the Fourier transform
\[
(\mathcal F\phi)(\xi)
 =\int_{\mathbb R}\phi(x)e^{-2\pi ix\xi}\,dx.
\]
It restricts to an isomorphism \(V\to V\): it preserves evenness
and the Schwartz space, and interchanges evaluation at zero with
the integral, so preserves the two zero moments in
\eqref{CFStage:BF:BF:eq:spaces}.  Its inverse on \(V\) is the inverse Fourier
transform with the same convention.
For the two-leg differential and its cycle map,
\[
d_{\rm leg}:V\oplus V\longrightarrow\mathcal B,\qquad
d_{\rm leg}(v,w)=\Theta(v-\mathcal Fw),
\]
\[
\gamma:V\longrightarrow V\oplus V,\qquad
\gamma(v)=(v,\mathcal F^{-1}v),
\]
one has \(d_{\rm leg}\gamma=0\) by direct substitution.
For a linear map \(L:E_Z\to V\), adding \(\gamma L\) to a
primitive therefore leaves \(B_n\)
unchanged, and hence leaves \(W_n\) unchanged for fixed \(R_n\).
The support transition described above likewise records where
the boundary becomes a supported zero while retaining the
pairing that produced \eqref{CFStage:BF:BF:eq:rank-two}.

There is also an exact tensor obstruction independent of the
rank-two argument.  For an integer \(k\ge1\), keep
\begin{equation}
A_k=\sum_{j=1}^k
 I^{\otimes(j-1)}\otimes A\otimes I^{\otimes(k-j)}
 \quad\hbox{on }E_Z^{\otimes k}.
\tag{BF.40}\label{CFStage:BF:BF:eq:tensor-generator}
\end{equation}
For any positive definite metric \(H\) on this same tensor
space, set \(W_{k,H}=A_k^*H+HA_k-kH\).
The nonzero tensor \(v^{\otimes k}\) of the eigenvector above
satisfies \(A_kv^{\otimes k}=k\rho v^{\otimes k}\).
Consequently
\begin{equation}
\frac{(v^{\otimes k})^*W_{k,H}v^{\otimes k}}
     {(v^{\otimes k})^*H v^{\otimes k}}
 =k(2\operatorname{Re}\rho-1).
\tag{BF.41}\label{CFStage:BF:BF:eq:tensor-floor}
\end{equation}
The denominator is strictly positive.
Thus the least two-sided excess for \(H,W_{k,H}\) is at least
\(k|2\operatorname{Re}\rho-1|\).
In particular this applies to the actual orbit metric of
A1721, lines 4102--4139,
\[
H_{k,T}
 =\int_0^T e^{-kt}\exp(tA_k)^*G_n^{\otimes k}
                       \exp(tA_k)\,dt,\qquad T>0.
\]
It is positive definite because the integrand has positive
quadratic form on every nonzero vector at every \(t\);
\(\exp(tA_k)\) is invertible.
The same time parameter is used in every factor, and the full
weight \(e^{-kt}\) remains.  Changing \(T\) cannot change the
Rayleigh value \eqref{CFStage:BF:BF:eq:tensor-floor}.

For the empty packet, \(h_Z=1\), \(E_Z=\{0\}\), and every
packet sum is zero.  Define its control excess to be zero.
Then the boundary-floor assertion is the equality \(0=0\).
The extension row \([t^{d-1}]R_Z\) and the positive-dimensional
Gram inverses above are not invoked in this empty case.

Equations \eqref{CFStage:BF:BF:eq:one-sided-mass} and
\eqref{CFStage:BF:BF:eq:boundary-floor} are lower bounds proved on the
actual arithmetic objects.  They neither assert an off-line
zero nor provide the arithmetic upper estimate needed to
exclude one.  The source explicitly retains that missing
estimate in A1721, line 4259, and A1771, line 4890.
Finite regression checks, the stated integer-sum tail bound,
and absolute Gram decrements do not supply it.

